Fraud schemes target drivers with injury claims that follow a traffic violation without contact and arrive days later; in 2026, Hong Kong police reported more than 70 such cases. Crash-triggered recorders store nothing; the accused vehicle cannot prove the claimant was elsewhere. We formulate this as proof of absence and propose an image-free evidence system. From the claim, a verifier derives a road-frame zone and a clock-bounded window, excludes sensors failing self-checks or liveness, and grades its verdict from two refutation tiers: a rule tier reading the detection log, and a formal tier reading a signed free-space record, disclosed only for the zone with a check that nothing was withheld, refuting only when the bound on missing the claimed vehicle is at most $10^{-3}$. In simulation, both tiers made no wrong verdict on 5,302 nominal genuine and empty claims and refuted 68\% (rule) and 54\% (formal) of fabricated claims. In a pre-registered test on 40,000 claims under unseen degradations, the tiers made similar numbers of wrong refutations (6 and 5 of 24,993) and could not be separated; refutations agreed by both were wrong for 0.020\% of genuine claims, although the 0.1\% criterion failed within two conditions. With an occluder outside the zone, only the formal tier abstained on unobserved space. Ablations and attacks credit the claim-derived zone, persistence, attestation, liveness and abstention on weak reflectors. The free-space record costs 2.45~MB per driving hour and about 49~kB per disclosed enriched claim; the detection log is not yet sealed.
